\documentclass[sigconf]{acmart}

\settopmatter{printacmref=false}
\renewcommand\footnotetextcopyrightpermission[1]{}
\pagestyle{plain}

\title{ControlPoint-AI: A Framework-Agnostic Type-3 Security Extension for Policy-Enforced AI Agent Interactions}

\author{Yeamin}
\affiliation{%
  \institution{IndustrySphere AI}
  \city{Dhaka}
  \country{Bangladesh}
}
\email{contact@industrysphere.ai}

\begin{abstract}
Autonomous AI agents are increasingly entrusted with access to external software tools, model APIs, and enterprise databases. However, relying on natural language system prompts or coarse-grained API keys to enforce security boundaries exposes organizations to indirect prompt injection, privilege escalation, unauthorized data exfiltration, and resource exhaustion. This paper introduces ControlPoint-AI, a lightweight, pluggable, framework-agnostic Type-3 security extension designed for the NDSS SafeOSE '27 Area 03 specification. Operating independently of internal neural network weights or hidden model reasoning, ControlPoint-AI intercepts structured agent-tool interactions at two deterministic control points: pre-execution and post-execution. ControlPoint-AI integrates natively with the Natural Language Interaction Protocol (NLIP) using the official nlip-sdk (v0.1.3) and maintains tamper-evident audit provenance via sequential SHA-256 hash chains. In an empirical evaluation across 15 synthetic enterprise attack scenarios, ControlPoint-AI achieved a 100.0\% attack prevention rate with a median P50 decision latency of 9.28 ms, comparing favorably against unprotected execution and conventional basic authentication.
\end{abstract}

\begin{document}
\maketitle

\section{Introduction}
Modern autonomous agent frameworks (e.g., LangChain, LlamaIndex, CrewAI) empower Large Language Models (LLMs) to formulate tool execution plans and invoke external APIs. While this paradigm unlocks workflow automation, it dissolves traditional security perimeters. If an autonomous agent encounters adversarial context, the model may generate tool calls that violate least-privilege principles, exfiltrate confidential data, or trigger destructive mutations.

To address these vulnerabilities, the NDSS SafeOSE '27 workshop established Area 03: Security Functions at Unified Control Points. Under this taxonomy, a Type-3 Security Extension governs interactions between the agent runtime and external resources using structured interaction data, standardized protocol messages, local state, and explicit administrative policies.

\section{Problem Definition and Threat Model}
Conventional agent deployments suffer from three structural security deficiencies:
\begin{enumerate}
    \item \textbf{Unchecked Tool Invocations}: Agents are often granted binary access to entire tools without field-level or operation-level constraints.
    \item \textbf{Observation Leakage}: Tools frequently return multi-field database rows containing sensitive attributes even when the agent requested only public attributes.
    \item \textbf{Audit Vulnerability}: Traditional flat logging systems cannot prove post-facto that audit entries have not been altered or deleted.
\end{enumerate}

\section{System Architecture}
ControlPoint-AI operates as an inline proxy between the autonomous agent runtime and external tool endpoints. It enforces four deterministic decisions: ALLOW, DENY, REDACT, and THROTTLE.
Pre-execution checks validate agent identity using salted SHA-256 tokens, sanitize argument syntax, check sliding-window rate limits, and evaluate explicit policies.
Post-execution inspection monitors tool observations and applies selective field-level redaction before returning data to the agent context stream.

\section{NLIP Protocol Integration}
Using the official nlip-sdk (v0.1.3), ControlPoint-AI processes incoming NLIP\_Message envelopes, translates them into canonical NormalizedRequest objects, and formats responses as compliant NLIP messages containing structured submessages with cryptographic audit proof tokens.

\section{Security Evaluation}
We evaluated ControlPoint-AI across 15 synthetic enterprise attack scenarios. ControlPoint-AI successfully resolved all 15 scenarios, achieving a 100.0\% attack prevention rate (APR) and 0.0\% false positive rate (FPR).

\begin{table}[h]
\caption{Comparative Tri-Baseline Evaluation}
\begin{tabular}{lccc}
\toprule
Metric & Baseline A (Unprotected) & Baseline B (Basic Auth) & Baseline C (ControlPoint-AI) \\
\midrule
Attack Success Rate (ASR) & 88.9\% & 44.4\% & \textbf{0.0\%} \\
Attack Prevention Rate (APR) & 0.0\% & 55.6\% & \textbf{100.0\%} \\
P50 Latency (ms) & 0.22 ms & 0.55 ms & \textbf{9.28 ms} \\
Throughput (RPS) & $\sim$4,500 & $\sim$1,850 & $\sim$890 \\
Audit Provenance & None & None & \textbf{SHA-256 Verified} \\
\bottomrule
\end{tabular}
\end{table}

\section{Conclusion}
ControlPoint-AI establishes verifiable interaction governance over autonomous AI agent tool calls. Complete source code, benchmark suites, and reproduction guides are open-source under Apache-2.0.

\end{document}
